\documentclass{leveldoc}
\DocUses{foundations,deflection_plates}
\DocFingerprint{80e3be3dabf8033a}

\begin{document}

\section{What the level asks}\label{sec:asks}

\leveldescription

One positive ion is launched along $+x$ from the left edge, on the mid-line
between two metal plates. Both plates are grounded until the player switches
on the supply attached to the upper plate and sets its potential. A negative
potential pulls a positive ion toward that plate; a positive potential pushes
the ion away. The lower plate stays grounded. The player sets no charges and
no magnets.

The ion must enter the detector. The shot has no acceptance window, no kinetic
window and no gate. A flight ends when the ion touches a plate, leaves the
world, or enters the detector (\xsecref{foundations}{sec:board}). A plate's
obstacle is its box plus a contact shell of thickness $\val{contact}$. The
ion's radius is $0$.

The reference switches the supply on and sets the upper plate to
$\val{vSupply}$. The ion climbs and meets the detector's left side. With the
supply off, both plates are grounded, the ion flies straight along the launch
line, and that line passes under the detector. This document uses the
deflection plates' rectangle collocation and uniform-gap limit, with the
supply entered as the upper plate's potential, and integrates the flight by
the classical Runge--Kutta method.

\section{Setup}\label{sec:setup}

\setuptable

The speed of light is $c = \val{cLight}$. The ion has charge $q = \val{qPart}$,
mass $m = \val{mPart}$ and radius $0$, and it is launched with kinetic energy
$T_0 = \val{T0}$. The launch is at $(\val{xLaunch}, \val{yLaunch})$, straight
along $+x$. Then $\gamma_0 = \val{gamma0}$, $\abs{\vect{p}_0} = \val{p0}$ and
$v_0 = \val{v0} = \val{vOverC}\,c$ (\xeqref{foundations}{eq:launch}).

The upper plate occupies $\val{xTop0} \le x \le \val{xTop1}$,
$\val{yTop0} \le y \le \val{yTop1}$. The supply sits on its centre
$(\val{xSupply}, \val{ySupply})$ and, in the reference, holds it at
$V_\mathrm{top} = \val{vTop}$. The lower plate occupies
$\val{xBot0} \le x \le \val{xBot1}$, $\val{yBot0} \le y \le \val{yBot1}$, at
$V_\mathrm{bot} = \val{vBot}$. The gap between the inner faces is
$g = \val{gap}$. The slider runs from $-\val{vSlider}$ to $\val{vSlider}$.

The detector is $\val{xDetLo} \le x \le \val{xDetHi}$,
$\val{yDetLo} \le y \le \val{yDetHi}$. The world is $-1 \le x \le \val{xWorld}$,
$-1 \le y \le \val{yWorld}$. The reported time step is $h = \val{H}$. The mesh
is the deflection plates' mesh: rectangles of edge $\val{panelFine}$,
$\val{nPanels}$ of them on the upper half.

\begin{figure}[ht]
\centering
\begin{tikzpicture}
\begin{axis}[axis equal image, width=0.92\linewidth,
  xmin=-1.5, xmax=32.5, ymin=-2, ymax=22.5,
  xlabel={$x$}, ylabel={$y$}, tick label style={font=\small}]
  \draw[black!40, dashed] (-1,-1) rectangle (\rawval{xWorld},\rawval{yWorld});
  \draw[black!25] (0,0) rectangle (30,20);
  \fill[green!50!black, fill opacity=0.18]
    (\rawval{xDetLo},\rawval{yDetLo}) rectangle (\rawval{xDetHi},\rawval{yDetHi});
  \draw[green!40!black]
    (\rawval{xDetLo},\rawval{yDetLo}) rectangle (\rawval{xDetHi},\rawval{yDetHi});
  \fill[black!20]
    (\rawval{xTop0},\rawval{yTop0}) rectangle (\rawval{xTop1},\rawval{yTop1});
  \draw[orange!80!black, thick]
    (\rawval{xTop0},\rawval{yTop0}) rectangle (\rawval{xTop1},\rawval{yTop1});
  \fill[black!20]
    (\rawval{xBot0},\rawval{yBot0}) rectangle (\rawval{xBot1},\rawval{yBot1});
  \draw[black!70]
    (\rawval{xBot0},\rawval{yBot0}) rectangle (\rawval{xBot1},\rawval{yBot1});
  \addplot[orange!80!black] table {pathbare.dat};
  \addplot[very thick, blue!60!black] table {pathref.dat};
  \draw[-{Stealth[length=5pt]}, thick] (\rawval{xLaunch},\rawval{yLaunch}) -- (\rawval{arrowX},\rawval{arrowY});
  \fill (\rawval{xLaunch},\rawval{yLaunch}) circle[radius=0.08];
\end{axis}
\end{tikzpicture}
\caption{The plates and the two flights. Both rectangles are the plates; the
upper one, outlined, is the plate the supply sets. The green box is the
detector. Blue is the reference flight: the negative upper plate pulls the ion
up, and the ion meets the left side of the detector. Orange is the same launch
with the supply off and both plates grounded: the ion stays on the launch line
and leaves through the right edge of the world, below the detector. Both paths
stay outside the plates.}
\label{fig:board}
\end{figure}

\section{Principles and model}\label{sec:model}

\begin{principles}
  \item The relativistic equation of motion \xeqref{foundations}{eq:eom}, the
    launch \xeqref{foundations}{eq:launch}, and energy in a static field
    \xeqref{foundations}{eq:energy}. There is no magnetic field, and the motion
    stays in the plane (\xsecref{foundations}{sec:plane}).
  \item A supply sets the potential of the tunable plate it sits on
    (\physref{2.7}). A grounded plate is potential $0$. The surface charge is
    the deflection plates' collocation
    \xeqref{deflection_plates}{eq:collocation}, with the rectangle integral
    \xeqref{deflection_plates}{eq:rectangle}.
  \item The uniform gap \xeqref{deflection_plates}{eq:uniform} is the analytic
    limit of that collocation. The flight uses the panels.
  \item The board's rules (\xsecref{foundations}{sec:board}) and the classical
    Runge--Kutta method with cubic Hermite event location
    (\xsecref{foundations}{sec:numerics}).
  \item The radiated power \xeqref{foundations}{eq:larmor}, summed by Simpson's
    rule \xeqref{foundations}{eq:simpson}. It is not put back into the force.
\end{principles}

The reference potential on the upper plate is the supply value, and the lower
plate is grounded. Off the metal the field is the field of the surface charge
and of its mirror. There is no placed charge. The state $(\vect{x}, \vect{p})$
obeys \xeqref{foundations}{eq:eom} with $\vect{B} = \vect{0}$, from the launch
until the first ending event.

\begin{assumption}{The supply is static}\label{as:static}
The potential equals the slider value at once. There is no resistor and no
ramp. The launch speed is $\val{vOverC}\,c$. The game uses the same static
supply (\physref{2.7}).
\end{assumption}

\begin{assumption}{No radiation reaction}\label{as:radiation}
The ion radiates, and that energy is not taken from the flight. On the
reference flight the radiated energy is a fraction $\val{RefRadFrac}$ of $T_0$.
\end{assumption}

\begin{assumption}{No image of the ion}\label{as:image}
The ion's own surface charge is left out, as the deflection plates leave it
out and as the game leaves it out (\physref{2.7}). The neglected force is
bounded by $q^2/d^2$, with the distance $d$ floored at $10^{-3}$, relative to
$\max(\abs{q\vect{E}}, T_0)$. The reference samples reach $\val{RefImage}$; the
grounded-plate samples reach $\val{BareImage}$. The samples are the step ends
and the ending point. The engine asks for the ratio to stay below $10^{-10}$.
\end{assumption}

\begin{assumption}{Equal rectangles}\label{as:panels}
The mesh is the one stated for the deflection plates: equal rectangles of edge
$\val{panelFine}$, not the engine's graded triangles. Replacing the edge by
$\val{panelCoarse}$ moves an end state by $\val{errMesh}$.
\end{assumption}

\begin{assumption}{The fixed step}\label{as:step}
The reported step is $h = \val{H}$. Halving it estimates an end state by
$\val{errStep}$. Ten times the panel-size estimate $\val{errMesh}$ lies above
$10^{-9}$, and it is larger than the step-halving estimate $\val{errStep}$, so
the tolerance on a state is that product, $\val{tolFlight}$. Between samples
the ion moves at most $\val{stepReach}$. The smallest sampled clearance is
$\val{RefDmin}$ on the reference flight and $\val{BareDmin}$ with the supply
off.
\end{assumption}

\section{Derivation}\label{sec:derivation}

\subsection{The supply and the uniform gap}\label{sec:gap}

The supply element at $(\val{xSupply}, \val{ySupply})$ is the player's setting
of the tunable electrode centred there. In the reference its value is
$\val{vSupply}$, so $V_\mathrm{top} = \val{vTop}$ and $V_\mathrm{bot} = \val{vBot}$.
With the supply omitted, both potentials are $0$.

The infinite-plate field is \xeqref{deflection_plates}{eq:uniform}. With these
two potentials it equals $\val{eUniform}$, pointing from the grounded plate up
toward the negative plate. Integrating that constant field from the lower
plate, the potential is linear, so at the midpoint of the gap it is the
average of the two voltages,
\begin{equation}\label{eq:mid}
  \phi_\mathrm{mid} = \frac{V_\mathrm{bot} + V_\mathrm{top}}{2} = \val{phiUniform}.
\end{equation}
These plates are finite. Halfway along them, on the launch line, the panel
field is $\val{eCenter}$ and the potential is $\val{phiCenter}$. The field is
lower than \xeqref{deflection_plates}{eq:uniform} by the relative amount
$\val{eUniformRel}$, and the potential differs from \eqref{eq:mid} by the
relative amount $\val{phiUniformRel}$. The march uses the panels. With both
plates grounded the panel field at that same point is $\val{eBare}$.

\subsection{The panels and the flight}\label{sec:flight}

The rectangles, the mirror and the collocation are
\xeqref{deflection_plates}{eq:rectangle} and
\xeqref{deflection_plates}{eq:collocation}. The right-hand side is the plate
potential, because there is no placed charge. The reported mesh has
$\val{nPanels}$ upper rectangles. The residual at the centroids is
$\val{solveRes}$. On the inner faces, off the centroids, the potential differs
from the plate potential by at most $\val{surfRes}$.

The conserved diagnostic is the deflection plates' flight energy
\xeqref{deflection_plates}{eq:flight-energy}: kinetic energy plus $q\phi$, with
no image term. The state is advanced by \xeqref{foundations}{eq:rk4} with step
$h = \val{H}$. Events are located as on the deflection plates: thirty-two
Hermite samples on the step, then sixty bisections, obstacles before the
world, and the world before the detector. A face coordinate is assigned when
the residual is below $10^{-6}$. The step-halving estimate divides the change
by $15$; the panel-size estimate is the change itself when the edge doubles.

\section{The reference solution by hand}\label{sec:hand}

\begin{procedure}
  \item\label{st:launch} \emph{The launch} \xeqref{foundations}{eq:launch}.
    $\gamma_0 = \val{gamma0}$, $\abs{\vect{p}_0} = \val{p0}$, so
    $\vect{p}_0 = (\val{p0}, 0)$ and $v_0 = \val{v0}$.
  \item\label{st:supply} \emph{The supply.} Set the upper plate to
    $\val{vSupply}$ and leave the lower plate at $\val{vBot}$. The uniform-gap
    limit is $E_y = \val{eUniform}$, and the linear mid-gap potential is
    $\val{phiUniform}$.
  \item\label{st:panels} \emph{The panels.} Solve
    \xeqref{deflection_plates}{eq:collocation} on the mesh of
    \S\ref{sec:flight}. The centroid residual is $\val{solveRes}$. The
    off-centroid discrepancy on the inner faces is $\val{surfRes}$. At the
    centre of the gap the panel field is $\val{eCenter}$ and the potential is
    $\val{phiCenter}$.
  \item\label{st:march} \emph{The march.} Advance
    \xeqref{deflection_plates}{eq:flight-energy} and
    \xeqref{foundations}{eq:eom} by \xeqref{foundations}{eq:rk4} with step
    $\val{H}$, and locate the ending face as in \S\ref{sec:flight}. The state
    every five time units:
    \begin{center}\small\val{states}\end{center}
  \item\label{st:end} \emph{The arrival.} The ion meets the \val{RefFace} face
    of the detector at $(\val{RefX}, \val{RefY})$, at $t = \val{RefT}$, with
    momentum $(\val{RefPx}, \val{RefPy})$, at an angle $\val{RefAng}$ degrees
    above $+x$. The abscissa is the detector's left edge, and the ordinate
    lies between $\val{yDetLo}$ and $\val{yDetHi}$. The run takes
    $\val{RefSteps}$ steps. The path only climbs: its lowest point is the
    launch, and its highest is $\val{RefYmax}$.
  \item\label{st:check} \emph{The checks.} Along the flight, $W$ drifts by at
    most $\val{RefWdev}$. The step-halving estimate is $\val{RefErr}$, and the
    larger of the two flights is $\val{errStep}$. The panel-size estimate is
    $\val{errMesh}$. The sampled clearance is $\val{RefDmin}$, outside the
    shell $\val{contact}$. The image-force ratio reaches $\val{RefImage}$.
  \item\label{st:rad} \emph{The radiation.} Simpson's rule on steps of $h/2$
    gives $\val{RefRad}$, a fraction $\val{RefRadFrac}$ of $T_0$. Halving that
    step again changes the integral by the relative amount $\val{RefRadErr}$.
    This integral is not a row of the comparison.
\end{procedure}

\begin{result}[The reference flight]
The ion enters the detector through its left side at
$(\val{RefX}, \val{RefY})$ at $t = \val{RefT}$, with momentum
$(\val{RefPx}, \val{RefPy})$. The level is solved.
\end{result}

\section{Comparison with the game}\label{sec:comparison}

\comparisontable

Eight comparisons are in the table: the reference outcome, time and end
position, and the outcome, time and end position with the supply off. The
tolerance on a state is $\val{tolFlight}$. Ten times the panel-size estimate
$\val{errMesh}$ lies above $10^{-9}$, and it is larger than the step-halving
estimate $\val{errStep}$, so the tolerance is that product. The tolerance on a
state was fixed from these estimates before the comparison.

The largest absolute difference is the reference ordinate,
$\val{cmpRefYDiff}$. It lies above the panel-size estimate $\val{errMesh}$ and
under the tolerance. The reference time differs by $\val{cmpRefTDiff}$, and
the bare time by $\val{cmpBareTDiff}$. The step estimate $\val{errStep}$ is far
below the reference ordinate: the panels, not the time step, set that
difference. The outcomes agree.

The reference abscissa, the bare abscissa and the bare ordinate agree exactly.
The reference abscissa is the detector's left edge, the bare abscissa is the
world's right edge, and the bare ordinate is the launch line.

\section{Other solutions}\label{sec:others}

The supply off, which is the empty placement in \code{placements.json}, does
not solve the level. Both plates are then grounded. The panel field at the
centre of the gap is $\val{eBare}$, and the ion stays on the launch line. It
leaves through the \val{BareFace} edge of the world at
$(\val{BareX}, \val{BareY})$, at $t = \val{BareT}$, with momentum
$(\val{BarePx}, \val{BarePy})$. The ordinate $\val{BareY}$ is the launch
height, and the detector begins at $\val{yDetLo}$, above that line. The
clearance of the plates is $\val{BareDmin}$, the energy drifts by
$\val{BareWdev}$, and the run takes $\val{BareSteps}$ steps.

The reference works because $V_\mathrm{top} = \val{vTop}$ is negative, so the
field in the gap points upward and the positive ion climbs from
$\val{yLaunch}$ to $\val{RefY}$, inside the detector's vertical span. The
slider's half-range is $\val{vSlider}$. This document integrates the reference
setting and the off setting, not the rest of the slider.

\section{What the game computes}\label{sec:engine}

A power supply sets the potential of the tunable electrode it sits on, and the
factorization of the panel system is reused when only that potential changes
(\physref{2.7}). The grounded lower plate is potential $0$. The surface charge
is the engine's triangle collocation, the same boundary condition as
\xeqref{deflection_plates}{eq:collocation}. This document uses the equal
rectangles of the deflection plates. The end states move by $\val{errMesh}$
when that edge doubles, and the reference ordinate in \S\ref{sec:comparison}
lies above that estimate and under ten times it.

The ion's image force is omitted on both sides. The bound sampled here reaches
$\val{RefImage}$, under the engine's limit $10^{-10}$. The flight in the
engine is the adaptive method of \physref{5} and \physref{6}, verified as in
\physref{7}. Our step estimate is $\val{errStep}$. Radiation is integrated and
not applied to the momentum (\xsecref{foundations}{sec:radiation},
\physref{8}).

\section{Further reading}\label{sec:reading}

The rectangle integral, the collocation and the uniform gap are derived for
the deflection plates (\xeqref{deflection_plates}{eq:rectangle},
\xeqref{deflection_plates}{eq:collocation},
\xeqref{deflection_plates}{eq:uniform}). The motion and the Runge--Kutta step
are in the Foundations (\xsecref{foundations}{sec:numerics}). The engine's
supplies and electrodes are \physref{2.7}.

\levelfileappendix

\end{document}
